\documentclass[11pt]{article}
\usepackage[margin=1in]{geometry}
\usepackage{hyperref}

\title{H2 Next Experiment Memo (2026)}
\author{C.M. Baird et al.}
\date{March 2026}

\begin{document}
\maketitle

\section*{Recommendation}
The single strongest next experimental move is a \textbf{Phase-0 benchtop optical visibility pilot} for H2. This is not the final decisive flagship platform; it is the shortest honest bridge from protocol/simulation to real hardware.

\section*{Why H2 first}
\begin{itemize}
    \item It is the cleanest physics-native falsifier in the package.
    \item It does not rely on QRNG trust assumptions.
    \item It maps directly to one measurable parameter, $\Gamma$.
    \item A null result is still scientifically valuable because it sets a real instrument-bounded floor.
\end{itemize}

\section*{Observable and decision quantity}
The target observable is interferometric visibility:
\[
\frac{V}{V_0} = \exp(-\Gamma T \Delta x^2).
\]
The key quantity to report is the achieved exclusion floor
\[
\Gamma_{\mathrm{floor}} = \frac{|\ln(1-\delta_{\mathrm{total}})|}{T\Delta x^2},
\]
where $\delta_{\mathrm{total}}$ is the full nuisance-budget precision floor.

\section*{Recommended first implementation}
\begin{itemize}
    \item Instrument class: benchtop optical interferometer (Mach-Zehnder or equivalent)
    \item Role: validate drift accounting, nuisance budgeting, and operator discipline
    \item Claim boundary: do not treat the pilot as a final flagship discovery test
\end{itemize}

\section*{Reference planning profile}
\begin{itemize}
    \item $T = 10^{-6}\,\mathrm{s}$
    \item $\Delta x = 10^{-3}\,\mathrm{m}$
    \item $\delta_{\mathrm{meas}} = 10^{-4}$
    \item timing fraction $=10^{-4}$
    \item path-separation fraction $=5\times10^{-4}$
    \item vibration fraction $=5\times10^{-4}$
    \item EM fraction $=10^{-4}$
    \item thermal fraction $=2\times10^{-4}$
\end{itemize}

Using the current reference budget, the derived pilot floor is approximately:
\[
\delta_{\mathrm{total}} \approx 1.15\times 10^{-3},
\qquad
\Gamma_{\mathrm{floor}} \approx 1.15\times 10^9\,\mathrm{s}^{-1}\mathrm{m}^{-2}.
\]

\section*{What success looks like}
\begin{itemize}
    \item Baseline-before and baseline-after remain stable within preregistered tolerance.
    \item Nuisance regressors do not explain away the inferred $\Gamma$.
    \item The run yields a transparent achieved $\Gamma_{\mathrm{floor}}$ under real instrument conditions.
\end{itemize}

\section*{Interpretation}
\begin{itemize}
    \item \textbf{Positive pilot:} escalate immediately to a more decisive flagship platform.
    \item \textbf{Null but stable pilot:} still a win, because it validates the real-world H2 chain.
    \item \textbf{Messy pilot:} also a win, because it reveals where the protocol is too fragile.
\end{itemize}

\section*{Bottom line}
If the question is ``what should we physically do next?'', the answer is: run the H2 benchtop optical pilot, publish the full nuisance budget and deviation log, and use that result to decide whether escalation to the flagship platform is warranted.

\end{document}
